\chapter{代码实践卷：Executable Evidence Kit}

本章对应仓库目录：

\begin{lstlisting}[language=bash]
books/cpu-volume-1-practice/labs/executable_evidence
\end{lstlisting}

这不是附带样例，而是第一册的主线实践工程。读者应该把它当成一条从零到研究报告的流水线：先让代码能解析一个人工构造的最小 ELF；再让它检查真实二进制；然后把 readelf、objdump、perf 和 benchmark 接进报告。每一步都保留测试和证据，不靠口头解释。

\section{一、构建、运行和目录结构}

构建命令如下：

\begin{lstlisting}[language=bash]
cmake -S books/cpu-volume-1-practice \
      -B books/cpu-volume-1-practice/build/executable-evidence-debug \
      -DCMAKE_BUILD_TYPE=Debug
cmake --build books/cpu-volume-1-practice/build/executable-evidence-debug
ctest --test-dir books/cpu-volume-1-practice/build/executable-evidence-debug \
      --output-on-failure
\end{lstlisting}

CLI 可以检查自己：

\begin{lstlisting}[language=bash]
books/cpu-volume-1-practice/build/executable-evidence-debug/labs/executable_evidence/cpu1ee_cli \
  books/cpu-volume-1-practice/build/executable-evidence-debug/labs/executable_evidence/cpu1ee_cli
\end{lstlisting}

目录结构按责任分层：

\begin{longtable}{p{0.32\linewidth}p{0.58\linewidth}}
路径 & 责任 \\
\hline
\code{include/cpu1/executable\_evidence.hpp} & 对外合同：输入配置、section 摘要、symbol 摘要、report 和格式化函数 \\
\code{include/cpu1/elf\_fixture.hpp} & synthetic ELF 构造器，供测试和 benchmark 共享 \\
\code{src/executable\_evidence.cpp} & ELF64 little-endian 解析、checksum、report 和 CSV 输出 \\
\code{src/main.cpp} & CLI 入口，只负责参数、文件读取和输出模式 \\
\code{tests/executable\_evidence\_tests.cpp} & GTest/GMock 测试，覆盖正常 fixture 和错误边界 \\
\code{bench/executable\_evidence\_bench.cpp} & Google Benchmark smoke，只测解析路径 \\
\end{longtable}

\section{二、分阶段实现和习题}

\topic{阶段 0：字节和 checksum。} 第一阶段不要急着解释 ELF。先把文件读成 \code{std::vector<std::uint8\_t>}，用 \code{checksum\_bytes} 得到稳定摘要。习题是构造两个只差一个字节的输入，证明 checksum 不同；再用同一输入重复计算，证明 checksum 稳定。

\topic{阶段 1：ELF magic 和 header。} 先只检查前四个 magic 字节，再检查输入至少有 ELF64 header 长度，最后读取 class、endianness、object type、machine 和 entry address。习题是写三个坏输入：完全不是 ELF、只有 magic 但太短、ELF32 class。每个坏输入都必须得到明确 diagnostic。

\topic{阶段 2：section table。} header 给出 section header offset、entry size、count 和 section-name string table index。实现时先读 raw section，再用 \code{.shstrtab} 把 name offset 翻译成名字。习题是解释为什么 section 名字不能在读 raw section 时立刻解析：因为当时还不知道 string table 的字节范围是否安全。

\topic{阶段 3：symbol table。} 符号表依赖两个事实：symbol table 自己的位置和大小，以及它 link 到的字符串表。习题是找到 \code{\_start} 的 symbol，说明 binding、type、section index、value 和 size 各自表达什么。不要把 value 直接叫“文件偏移”，它通常是虚拟地址或链接视角的地址。

\topic{阶段 4：CLI 和 CSV。} CLI 的职责是薄的：读取文件，调用 library，选择 report 或 CSV。习题是用 \code{--sections-csv} 输出 section 表，再写脚本检查表头是否固定为 \code{name,type,flags,address,offset,size}。

\topic{阶段 5：benchmark smoke。} Google Benchmark 只测 synthetic ELF 的解析路径，不测文件读取和打印。习题是解释为什么输入构造必须放在计时循环外：否则测到的是构造成本、分配成本和缓存状态，不是解析器本身。

\section{三、每章实践任务矩阵}

\begin{longtable}{p{0.24\linewidth}p{0.34\linewidth}p{0.34\linewidth}}
第一册章节 & 实践任务 & 验收证据 \\
\hline
心智模型 & 画出 source、compiler、ELF、loader、process 的证据链 & 报告能区分文件事实和运行事实 \\
程序到进程 & 用 CLI 检查一个自己编译的小程序 & 保存 byte size、checksum、machine、entry \\
编译器 IR 与优化 & 分别检查 Debug 和 Release 产物 & 说明 section、symbol 或 size 差异 \\
CPU 执行指令 & 对一个函数用 objdump 找机器码 & 能指出函数符号和 \code{.text} 的关系 \\
数据、内存、指针、布局 & 解释 section offset、size 和字节范围检查 & 不越界读取，不把虚拟地址当文件偏移 \\
Linux 工具与调试 & 用 readelf 交叉验证三个字段 & 字段名、数值和工具命令写入报告 \\
汇编语法 & 对照反汇编中的操作数和地址 & 能说明至少一条指令访问寄存器还是内存 \\
整数指令 & 检查解析器的 little-endian 读取函数 & 有 u16/u32/u64 小端样例测试 \\
控制流 & 解释错误输入为什么提前返回 & diagnostics 覆盖 magic、短 header、不支持 class \\
System V ABI & 找一个未被 inline 的函数符号 & 说明函数名、binding、type、section index \\
C++ 与汇编互操作 & 比较 C++ 名字、符号名字和工具输出 & 写清 name mangling 或优化影响 \\
测量基础 & 运行 benchmark smoke & 说明只测解析，不测 I/O 和打印 \\
benchmark 设计 & 写一份坏测量分析 & 指出计时边界里混入的无关成本 \\
\end{longtable}

\section{四、递进研究项目}

研究项目分六次提交，不要求一次写完。

\begin{longtable}{p{0.16\linewidth}p{0.32\linewidth}p{0.42\linewidth}}
阶段 & 交付物 & 通过标准 \\
\hline
0 & synthetic ELF 测试 & 正常 fixture 能解析出 \code{.text} 和 \code{\_start} \\
1 & 错误输入测试 & 非 ELF、短 header、不支持 class 都有 diagnostic \\
2 & 真实二进制 report & 对一个自己编译的程序输出 checksum、machine、entry、sections \\
3 & 工具交叉验证 & readelf 或 objdump 至少验证三个字段 \\
4 & 符号与 ABI 观察 & 选择一个函数解释 symbol、section、地址和优化影响 \\
5 & benchmark 与测量限制 & benchmark smoke 通过，报告说明未测 I/O、打印和真实 workload \\
6 & 最终研究报告 & 结论只来自测试、CLI 输出、工具输出和明确限制 \\
\end{longtable}

\begin{labbox}
最终报告的核心不是“我看到了很多字段”，而是“这些字段怎样互相证明”。例如：\code{cpu1ee\_cli} 输出 machine 为 x86-64，readelf header 同样显示 x86-64；\code{.text} 的 offset 和 size 能解释 objdump 反汇编来自哪里；symbol table 中的函数名能和源码函数对应；checksum 证明报告绑定的是这一次构建产物。
\end{labbox}
